% Part: computability
% Chapter: computability-theory
% Section: non-comp-set

\documentclass[../../../include/open-logic-section]{subfiles}

\begin{document}

\olfileid{cmp}{thy}{ncp}
\olsection{Ana Himpunan kang Ora Komputabel}


Ing ndhuwur wis kita deleng yen saben himpunan
komputabel bisa dienumerasi sacara komputabel.
Apa kosok baline uga bener? Ing ngisor iki
dituduhake yen sacara umum ora bener.

\begin{thm}\ollabel{thm:K-0}
Tetepna $K_0$ minangka himpunan
$\Setabs{\tuple{e, x}}{\cfind{e}(x) \fdefined}$.
Mula $K_0$ bisa dienumerasi sacara komputabel
nanging ora komputabel.
\end{thm}

\begin{proof}
Kanggo ndeleng yen $K_0$ bisa dienumerasi
sacara komputabel, gatekna yen himpunan iku
domain fungsi~$f$ kang ditegesi minangka
\[
f(z) = \umin{y}{(\len{z} = 2 \land T((z)_0, (z)_1, y))}.
\]
Sebab, yen $\cfind{e}(x)$ ditegesi,
$f(\tuple{e, x})$ nemokake runtutan
komputasi kang mandheg; yen
$\cfind{e}(x)$ ora ditegesi, semono uga
$f(\tuple{e, x})$; lan yen $z$ malah ora
ngode pasangan, $f(z)$ uga ora ditegesi.

Kasunyatan yen $K_0$ ora komputabel yaiku
ora bisane masalah mandheg diputusake,
\olref[hlt]{thm:halting-problem}.
\end{proof}

Himpunan $K_0$ yaiku himpunan pasangan
$\tuple{e,x}$ kang ndadekake
$\cfind{e}(x) \fdefined$, yaiku
$\tuple{e,x} \in K_0$ yen lan mung yen
$\cfind{e}$ ditegesi (mandheg) ing
input~$x$. Mula himpunan iki uga diarani
``himpunan mandheg''. Himpunan
$K = \Setabs{e}{\cfind{e}(e) \fdefined}$
yaiku ``himpunan mandheg ing indekse dhewe''.
Himpunan iki kerep digunakake minangka
conto baku himpunan kang ora bisa diputusake.

\begin{thm}\ollabel{thm:K}
Himpunan mandheg ing indekse dhewe
$K = \Setabs{e}{\cfind{e}(e) \fdefined}$
iku !!{c.e.} nanging ora bisa diputusake.
\end{thm}

\begin{proof}
  Upamane $K$ bisa diputusake, yaiku fungsi
  karakteristike $\Char{K}$ komputabel. Tetepna
  \[d(e) = \begin{cases}
  1 & \text{yen\/ $\Char{K}(e) = 0$}\\
  \fundefined & \text{yen ora.}
  \end{cases}
  \] 
  Tetepna $k$ minangka indeks~$d$, yaiku
  $d \simeq \cfind{k}$. Mula $d(k)
  \simeq \cfind{k}(k)$. Iki nalisir
  kasunyatan yen $d(k) \fdefined$ yen lan
  mung yen $\cfind{k}(k) \fundefined$,
  kang manut saka definisi~$d$.

  $K$ minangka domain $f(x) = \umin{y}{T(x,x,y)}$
  mula uga !!{c.e.}
\end{proof}

\end{document}
